\section{Specialization of an actual algebraic class}
\label{spec:section}

We first establish a specialization statement for projective analytic
families.  This permits the Kuga--Satake construction to be made on a
marked period polydisc, where all its rational data are constant.  Throughout
this section, ``very general'' in a connected complex manifold means
outside a countable union of proper closed analytic subsets.

\subsection{Hilbert parameters over an analytic base}

\begin{lemma}[An analytic classifying map]
\label{spec:hilbert-map}
Let $p:\mathcal X\to D$ be a smooth proper holomorphic family of
nonempty, pure $n$-dimensional varieties over a connected complex manifold.
Suppose that an embedding
\[
 \mathcal X\hookrightarrow D\times\mathbb P^N
\]
over $D$ has been given.  There are a reduced complex quasi-projective
scheme $H$, a smooth projective algebraic family
$\mathcal U\to H$, and a holomorphic map
$\eta:D\to H^{\mathrm{an}}$ such that $\mathcal X$ is the embedded
analytic pullback of $\mathcal U$.
\end{lemma}

\begin{proof}
The embedded fibers have a common Hilbert polynomial $P$.
Indeed, Riemann--Roch expresses
$\chi(\mathcal X_t,\mathcal O_{\mathcal X_t}(k))$ in terms of the
fiber integrals of the Chern classes of the relative tangent bundle and
the hyperplane bundle.  These integrals are locally constant in a smooth
proper family, and $D$ is connected.

Let $\operatorname{Hilb}^{P}(\mathbb P^N)$ be the algebraic Hilbert scheme.
Fixed-polynomial boundedness supplies one sufficiently large integer $k$
for which every member $Y$ has
\[
 \dim \operatorname{im}\bigl(H^0(\mathbb P^N,\mathcal O(k))
       \longrightarrow H^0(Y,\mathcal O_Y(k))\bigr)=P(k),
\]
and its degree-$k$ ideal space defines the Hilbert scheme's closed
embedding in the corresponding Grassmannian.  These uniform statements
precede the construction of $\eta$; see
\cite[Theorem~3.2 and Proposition~3.8]{GrothendieckHilbert}.

We check that $t\mapsto H^0(\mathbb P^N,\mathcal I_{\mathcal X_t}(k))$
is holomorphic.  Fix $t_0\in D$.  Since $\mathcal X_{t_0}$ is reduced,
point evaluations separate the restrictions of degree-$k$ polynomials.
Choose finitely many points whose evaluations have rank $P(k)$, and
extend these points to local holomorphic sections of $p$, using its
submersion property.  After locally trivializing $\mathcal O(k)$ along
these sections, evaluation is a holomorphic matrix.  A nonzero
$P(k)$-minor remains nonzero near $t_0$.  Its rank cannot exceed $P(k)$,
because evaluation factors through restriction to $\mathcal X_t$.
Consequently its kernel equals the degree-$k$ ideal space, and these
kernels give the required holomorphic Grassmannian map.

Pointwise this map lies in the Hilbert locus.  Since $D$ is reduced,
the defining equations of that locus vanish on $D$, so the map factors
holomorphically through the Hilbert scheme.  Take $H$ to be the reduction
of the open locus parametrizing smooth pure $n$-dimensional subschemes,
and let $\mathcal U$ be the restricted universal family.  Flatness of the
universal family and smoothness of all its fibers imply that
$\mathcal U\to H$ is smooth.  Its pullback to $D$ and $\mathcal X$
are reduced closed analytic subspaces of $D\times\mathbb P^N$ with the
same fibers, hence the same underlying set.  They therefore coincide.
\end{proof}

For a closed subscheme $Z$ of dimension at most $a$, write
$[Z]_a$ for its $a$-dimensional fundamental cycle, with the lengths at
the generic points as multiplicities.  Components of smaller dimension
do not enter $[Z]_a$.

\begin{lemma}[Classes of flat subschemes]
\label{spec:flat-cycle}
Let $f:M\to R$ be a smooth projective algebraic family of relative
dimension $n$, with $R$ a smooth irreducible complex variety of dimension
$b$.  Let $Z\subset M$ be a closed subscheme, flat over $R$, whose fibers
have dimension $n-2$.  Then
\[
 r\longmapsto
 \operatorname{cl}([Z_r]_{n-2})\in H^4(M_r,\mathbb Q)
\]
is a section of the local system $R^4f^{\mathrm{an}}_*\mathbb Q$.
This includes nonpure and nonreduced fibers.
\end{lemma}

\begin{proof}
The total space $M$ is smooth.  Flatness and the dimension formula give
$\dim Z=n+b-2$, so its top cycle has a class
\[
 \zeta=\operatorname{cl}([Z]_{n+b-2})\in H^4(M,\mathbb Q).
\]
Here the usual cycle class on a smooth variety takes values in ordinary
cohomology; completeness of $M$ is unnecessary.
We claim that, for the inclusion $i_r:M_r\hookrightarrow M$,
\begin{equation}
 \label{spec:restriction}
 i_r^*\zeta=\operatorname{cl}([Z_r]_{n-2}).
\end{equation}

Choose local algebraic coordinates by an \emph{\'etale} map from
a neighborhood $U\subset R$ to $\mathbb A^b$ taking $r$ to $0$.
Shrink $U$ to remove the other points over $0$. The coordinate
functions then cut out a smooth flag
\[
 U=U_0\supset U_1\supset\cdots\supset U_b=\{r\}.
\]
Put $Z_j=Z\times_R U_j$. It remains flat over $U_j$ and has dimension
at most $n+b-2-j$. At each step the next coordinate is a
nonzerodivisor on $U_j$ and hence on the entire scheme $Z_j$.
The Cartier intersection formula therefore gives
\[
 [Z_j]_{n+b-2-j}\longmapsto[Z_{j+1}]_{n+b-3-j},
\]
with the scheme multiplicities
\cite[Lemma~42.29.5, Tag~02TB]{StacksCartier}.
Applying the formula successively yields
\[
 i_r^![Z]_{n+b-2}=[Z_r]_{n-2}.
\]
The formula applies to the whole possibly nonpure or nonreduced
scheme at each cut; no componentwise flatness assertion is needed.
Compatibility of Betti cycle classes with these regular pullbacks proves
\eqref{spec:restriction}.

Finally, the restrictions of a single cohomology class on the total
space of a smooth proper family form a section of its cohomology local
system.  This follows directly from local differentiable triviality.
Apply it to $\zeta$.
\end{proof}

\begin{proposition}[Specialization from very general algebraicity]
\label{spec:specialization}
Let $D$ be a connected polydisc, and let
$p:\mathcal X\to D$ be a smooth proper holomorphic family of
$n$-dimensional varieties, $n\geq2$, with a relative embedding in
$D\times\mathbb P^N$.  Let
\[
 \theta\in H^0(D,R^4p_*\mathbb Q)
\]
be a flat rational class.  If $\theta_t$ is the class of a rational
codimension-two algebraic cycle for very general $t\in D$, then the same
holds for every $t\in D$.
\end{proposition}

\begin{proof}
Fix a desired point $0\in D$ and trivialize $R^4p_*\mathbb Q$ over $D$.
Use Lemma~\ref{spec:hilbert-map} to write the family as the pullback of
$\mathcal U\to H$ along $\eta:D\to H^{\mathrm{an}}$.
For each possible Hilbert polynomial $Q$ of degree $n-2$, consider
\[
 T_Q=\operatorname{Hilb}^{Q}(\mathcal U/H).
\]
This is projective over $H$ and has a flat universal subscheme.
It has finitely many reduced irreducible components.  Resolve each
component by a projective birational morphism, obtaining a finite
collection of smooth irreducible varieties $R\to (T_Q)_{\mathrm{red}}$
whose images cover $(T_Q)_{\mathrm{red}}$; use
\cite{Hironaka}.  Thus each composite $R\to H$ is proper.
Pullback gives a smooth projective ambient family $\mathcal U_R\to R$
and a flat closed subscheme $Z_R\subset\mathcal U_R$.  By
Lemma~\ref{spec:flat-cycle}, the classes of the top cycles of its fibers
are sections of $R^4(\mathcal U_R/R)_*\mathbb Q$.

Form the analytic spaces
\[
 K_R=D\times_{H^{\mathrm{an}}}R^{\mathrm{an}}.
\]
They are proper over $D$.  Their pulled-back ambient families agree with
the pullback of $\mathcal X$, so the preceding cycle classes become
locally constant functions on $K_R$ with values in the fixed marked
space $H^4(\mathcal X_0,\mathbb Q)$.  The spaces $K_R$ may be singular
or nonreduced: this causes no difficulty, since local constancy was
established on the smooth parameter $R$ before base change.

There are countably many $Q$ and finitely many choices of $R$ for each
$Q$.  Each $K_R$ is second countable and has a locally finite family of
closed irreducible analytic components.  It consequently has at most
countably many such components.  Every component is proper over $D$;
Remmert's Proper Mapping Theorem makes its image a closed analytic
subset of $D$ \cite[\S7, Satz~1]{Grauert}.  Let $E$ be the union of the images
of all components whose image is not all of $D$.  Thus $E$ is a
countable union of proper closed analytic subsets.

Choose $t\in D$ outside $E$ and outside the exceptional sets in the
hypothesis.  Such a point exists by the Baire Category Theorem.  Choose
an actual cycle representing $\theta_t$, and write it as a finite sum
\[
 \gamma_t=\sum_{j=1}^{s}a_j[Y_{j,t}],
 \qquad a_j\in\mathbb Q,
\]
where each $Y_{j,t}\subset\mathcal X_t$ is an integral subvariety of
codimension two.  Its Hilbert point lifts to one of the parameters $R$,
and hence to a point $u_j\in K_R$ above $t$.  Choose an irreducible
component $C_j$ containing $u_j$.  Because $t\notin E$, the map
$C_j\to D$ is surjective.  In particular, there is a point
$v_j\in C_j$ above $0$.

An irreducible complex analytic space is connected.  The marked cycle
class is therefore constant on $C_j$, and the fiber of the universal
subscheme at $v_j$ has top cycle $\delta_{j,0}$ satisfying
\[
 \operatorname{cl}(\delta_{j,0})
 =\operatorname{cl}([Y_{j,t}])
 \quad\text{under the marking.}
\]
These are algebraic cycles on $\mathcal X_0$: the fibers of the
universal subscheme are algebraic in $\mathbb P^N$, and the embedded
fiber identification is algebraic by GAGA \cite{SerreGAGA}.
Consequently
\[
 \gamma_0=\sum_{j=1}^{s}a_j\delta_{j,0}
 \quad\text{satisfies}\quad
 \operatorname{cl}(\gamma_0)=\theta_0.
\]
The argument applies to any chosen point of $D$.
\end{proof}

The avoidance in this proof is essential: density of fibers with
algebraic classes by itself would not ensure that the components
parametrizing those cycles reach the desired fiber.  Proposition~\ref{spec:specialization}
requires actual algebraicity at very general points.  It assumes no
principle asserting algebraicity from the Hodge property of a flat
class.

\subsection{A marked family for the primitive Clifford tensor}

\begin{lemma}[Local polarized families]
\label{spec:local}
Fix a primitively amply polarized K3 surface $(S_0,h)$.
After choosing a sufficiently small full-dimensional polarized period
polydisc $D$ about its period, there are smooth proper holomorphic
families
\[
 \mathcal S\longrightarrow D,
 \qquad \mathcal A\longrightarrow D,
\]
with the following properties.
The K3 family has central fiber $(S_0,h)$, a relative ample line bundle,
and a marking of its integral second cohomology.
The primitive rational local system is a fixed quadratic space
$(V,q)$ of dimension $21$.
The abelian family has marked rational first cohomology
$W=C^+(V,q)$ with the full Clifford Hodge structures.
Fixed rational polarization and parity-identification data give a flat
family of standard tensor maps
\[
 \mu_t:V\longrightarrow H^1(\mathcal A_t,\mathbb Q)
                       \otimes H^1(\mathcal A_t,\mathbb Q).
\]
Finally, after shrinking $D$, the product
\[
 \mathcal X=\mathcal S\times_D\mathcal A\times_D\mathcal A
\]
admits a relative embedding in $D\times\mathbb P^N$.
\end{lemma}

\begin{proof}
Polarized deformation theory and local Torelli give the smooth K3
family over a $19$-dimensional polydisc on which the period map is a
local isomorphism; see \cite[Chapter~6, Section~2]{Huybrechts}.
Since $D$ is contractible, its cohomology local systems are marked.
The class $h$ remains of type $(1,1)$.  We recall why it is represented
by a line bundle on the total family.
Local differentiable triviality over the contractible base extends
$h$ to $H^2(\mathcal S,\mathbb Z)$.  The polydisc is Stein, and coherent
Leray and Cartan's Theorem~B identify
\[
 H^2(\mathcal S,\mathcal O_{\mathcal S})
 =H^0(D,R^2p_*\mathcal O_{\mathcal S}),
 \qquad p:\mathcal S\to D.
\]
The direct image on the right is locally free and commutes with base
change.  The image of $h$ in it vanishes fiberwise, hence vanishes.
The exponential sequence now gives a holomorphic line bundle
$\mathcal L$ with first Chern class $h$.

A high power of $\mathcal L$ is very ample on $S_0$ and has vanishing
higher cohomology there.  Proper cohomology and base change extend a
basis of its sections over a neighborhood of $0$.  They give a
relative embedding after shrinking: basepoint freeness and immersion
persist near the central fiber, and properness excludes pairs of
distinct points with the same image.  For the latter assertion, any
sequence of such pairs approaching the central fiber either limits
to distinct points, contrary to its embedding, or to a single point,
contrary to the local embedding furnished by the immersion.
This proves local relative projectivity, and also that $h$ stays ample.
The analytic tools used here are the proper direct-image and base-change
theorems and the Stein vanishing theorem; see
\cite[\S6, Hauptsatz~I; \S7, S\"atze~3,5; \S2, Satz~3]{Grauert}.

On the fixed Clifford space $W$, the Hodge filtration and rational
alternating polarization are those constructed in
Proposition~\ref{cliff:construction}.  To recall the holomorphic dependence,
if $\omega_t$ spans the period line, then
\[
 F^1W_{\mathbb C}
 =\operatorname{im}\bigl(L_{\omega_t}:C^-(V_{\mathbb C},q)
                                      \longrightarrow C^+(V_{\mathbb C},q)\bigr).
\]
This has constant half dimension and varies holomorphically.
The polarization can be fixed on a sufficiently small chart.
Choose a rational symplectic basis and let its integral span be a
lattice $\Lambda$ in $W$. It is unimodular for the chosen symplectic form and
does not change the rational Hodge structure.
Since $W$ is cohomology, the associated torus is explicitly
\[
 \mathcal A_t=(F^1W_t)^\vee/\operatorname{image}(\Lambda^\vee),
\]
where a lattice functional is restricted to $F^1W_t$.
Thus $H_1(\mathcal A_t,\mathbb Z)=\Lambda^\vee$,
$H^1(\mathcal A_t,\mathbb Z)=\Lambda$, and
$H^{1,0}(\mathcal A_t)=F^1W_t$, compatibly with the markings.
The polarization supplies the Riemann form on the dual lattice.
Equivalently, the marked filtrations give a holomorphic map to Siegel
space and $\mathcal A\to D$ is the pullback of its marked torus family.
The full markings supply, in particular, level-$3$ data. The spin
action is defined on the fixed Clifford space over the polydisc,
so there is no monodromy lifting problem on this base.

For completeness, in symplectic period coordinates this torus family
is
\[
 \mathbb C^g/(\mathbb Z^g+\tau(t)\mathbb Z^g),
 \qquad \tau(t)^{\mathsf t}=\tau(t),\quad
          \operatorname{Im}\tau(t)>0.
\]
The factors
\[
 \exp\bigl(-\pi i\,n^{\mathsf t}\tau(t)n
               -2\pi i\,n^{\mathsf t}z\bigr),
 \qquad m+\tau(t)n\in\mathbb Z^g+\tau(t)\mathbb Z^g,
\]
define a holomorphic theta line bundle, ample on every fiber.
The preceding extension-of-sections argument gives a relative
projective embedding after shrinking $D$.

Fix a rational nonisotropic vector $z_0\in V$ and the
polarization identification $j:W^\vee(-1)\xrightarrow{\sim}W$.
The maps
\[
 V(1)\longrightarrow\operatorname{End}(W),
 \qquad x\longmapsto\bigl(w\longmapsto xwz_0\bigr)
\]
give $\mu_t$ after untwisting and applying $j$ to the dual slot.
Clifford multiplication, $z_0$, and $j$ are constant rational data in
the markings.  Thus $\mu_t$ is flat as a rational tensor and is a Hodge
map in every fiber.  A Segre embedding combines the relative embeddings
of $\mathcal S$ and the two copies of $\mathcal A$.
\end{proof}

\begin{corollary}[The primitive map at every polarized fiber]
\label{spec:primitive}
Fix a polarized moduli component, and assume that at very general
points the nonzero algebraic input of Theorem~\ref{main:generic} holds.
Then at every $(S_0,h)$ in that component the standard full primitive
Clifford tensor is induced by a rational codimension-two algebraic
cycle.  The same conclusion holds if the hypothesis is made outside
countably many proper closed analytic subsets in a full-dimensional
polarized period chart about each desired point.
\end{corollary}

\begin{proof}
Take the families in Lemma~\ref{spec:local} about $(S_0,h)$.
Under the markings, extend $\mu_t$ by zero on $\mathbb Qh$ and on the
other degrees of $H^*(\mathcal S_t,\mathbb Q)$.
Poincar\'e duality and the K\"unneth decomposition give a unique flat
class
\[
 \theta\in H^0(D,R^4(\mathcal X/D)_*\mathbb Q)
\]
with this action.  More explicitly, for a rational basis $(v_i)$ of
$V$, let $(q^{ij})$ be the inverse cup-pairing matrix.  Then
\begin{equation}
 \label{spec:theta}
 \theta_t=\sum_{i,j}q^{ij}v_i\otimes\mu_t(v_j)
 \ \in\ V\otimes
 \bigl(H^1(\mathcal A_t)\otimes H^1(\mathcal A_t)\bigr)
 \ \subset\ H^4(\mathcal X_t,\mathbb Q).
\end{equation}

By Proposition~\ref{gen:genericity}, after excluding countably many
proper closed analytic subsets we have
$\operatorname{Hg}(V_t)=\operatorname{SO}(V,q)$.
At a point satisfying this condition and the assumed input,
Theorem~\ref{gen:generation} produces a cycle inducing $\mu_t$ on $V_t$.
Its full class can be made exactly $\theta_t$.
Indeed, on $\mathcal S_t$ the algebraic correspondence
\begin{equation}
 \label{spec:primitive-projector}
 \pi_{V,t}=[\Delta_{\mathcal S_t}]
 -[o_t]\times\mathcal S_t-\mathcal S_t\times[o_t]
 -\frac{h\times h}{\int_{\mathcal S_t}h^2},
 \qquad o_t\in\mathcal S_t,
\end{equation}
acts identically on $V_t$ and kills $H^0$, $H^4$, and $\mathbb Qh$.
There is no odd K3 cohomology.  Precompose the chosen cycle with
$\pi_{V,t}$ and postcompose with the algebraic projector onto
$H^1(\mathcal A_t)\otimes H^1(\mathcal A_t)$, as in
Lemma~\ref{cyc:projectors}.  These operations preserve codimension two.
The resulting cycle has the prescribed action in every source degree,
so Poincar\'e duality and K\"unneth identify its entire class with
\eqref{spec:theta}.  The point $o_t$ is needed only on this fiber; its
cohomology class, and hence the projector's cohomological action, is
canonical.

If the hypothesis is stated in algebraic moduli, its exceptional loci
pull back to proper closed analytic subsets of $D$.  They remain proper
because the full-dimensional polarized period chart has open image
locally in that component, up to a finite automorphism quotient.
Thus $\theta_t$ is algebraic very generally on $D$ in either version
of the hypothesis.  Proposition~\ref{spec:specialization} produces an
actual rational codimension-two cycle with class $\theta_0$.
Its action on $V$ is exactly $\mu_0$. Undoing the fixed invertible right multiplier on the output Hom slot, using Lemma~\ref{cliff:slots}, also algebraically realizes $m^q$ itself.
\end{proof}

Neither the original witness abelian varieties nor their cycles have
to form a family.  Theorem~\ref{gen:generation} first transfers the
input to the fixed local full Clifford family.  The specialization
argument then uses a finite representative of one sufficiently general
cycle, keeping its rational coefficients unchanged.
